% !TEX root = ../../../main.tex
% 所属：第一部分：单变量分析学 / 第7章：微分法
% 节：7.1 导数、微分与连续性
% 正文按原文件顺序保存；由 ../chapter7.tex 调用。
\section{导数、微分与连续性}
\begin{definition}
	$y=f(x)$ 在 $x_0$ 的某个邻域 $O\left(x_0, \rho\right)$ 上有定义，若 $\forall \Delta x\left(x_0+\Delta x \in O\left(x_0, \rho\right)\right)$ ，极限
	
	$$
	\lim _{\Delta x \rightarrow 0} \frac{f\left(x_0+\Delta x\right)-f\left(x_0\right)}{\Delta x}
	$$
	
	
	存在且有限，则称 $f(x)$ 在 $x=x_0$ 可导，该极限称为 $f(x)$ 在 $x=x_0$ 处的导数，记为 $f^{\prime}\left(x_0\right)$ .
	

\end{definition}	等价地，$f^{\prime}\left(x_0\right)$ 可以写作

$$
f^{\prime}\left(x_0\right)=\lim _{t \rightarrow x_0} \frac{f(t)-f\left(x_0\right)}{t-x_0}
$$


若 $f(x)$ 在 $(a, b)$ 上任一点处可导，则称 $f(x)$ 在 $(a, b)$ 上可导，记为$f\in D(a,b)$,若$f$的导函数也是连续的，则记为$f\in C^1(a,b)$
\begin{definition}
	$ y=f(x)$ 在 $O\left(x_0, \rho\right)$ 上有定义，$\forall \Delta x\left(x_0+\Delta x \in O\left(x_0, \rho\right)\right).$ 
	极限
	
	$$
	\lim _{\Delta x \rightarrow 0^{+}} \frac{f\left(x_0+\Delta x\right)-f\left(x_0\right)}{\Delta x}
	$$
	
	
	存在且有限，则称 $f(x)$ 在 $x=x_0$ 处右可导，记为 $f^{\prime}_{+}\left(x_0\right)$ ，称右导数.
	
	类似地，定义 $$f^{\prime}_{-}\left(x_0\right)=\lim _{\Delta x \rightarrow 0^{-}} \frac{f\left(x_0+\Delta x\right)-f\left(x_0\right)}{\Delta x}$$
	显然地，$f^{\prime}\left(x_0\right)$ 存在的充要条件是 $f_{+}^{\prime}\left(x_0\right)=f_{-}^{\prime}\left(x_0\right)=f\left(x_0\right)$
	
	若 $f^{\prime}_{+}\left(x_0\right) \neq f^{\prime}\left(x_0\right)$ ，或其中有不存在者，则称 $x_0$ 为 $f(x)$ 的不可导点
	
	与连续性类似的，我们可以给出一个函数在闭区间上可导的定义：$f\in D(a,b)$，且$f'_{+}(a),f'_{-}(b)$存在且有限
\end{definition}
\begin{definition}
	$y=f(x)$ 在 $x_0$ 的某个邻域 $O\left(x_0, \rho\right)$ 上有定义，若 $\Delta x \rightarrow 0$ 时，$$\quad \Delta y=f\left(x_0+\Delta x\right)-f\left(x_0\right)=A\left(x_0\right) \Delta x+o(\Delta x)$$则称 $y=f(x)$ 在 $x_0$ 处可微，记其微分为
	$$\mathrm{d} y|_{y=y_0}=A\left(x_0\right) \mathrm{d} x$$
\end{definition}
\begin{theorem}
	$y=f(x)$在$x=x_0$处可导的充要条件是$f(x)$在$x_0$处可微.
\end{theorem}
\begin{proof}
	必要性：已知 $f(x)$ 在 $x_0$ 处可导，由定义
	$$f^{\prime}\left(x_0\right)=\lim _{\Delta x \rightarrow 0} \frac{f\left(x_0+\Delta x\right)-f\left(x_0\right)}{\Delta x}$$
	存在一个函数 $\varphi(\Delta x) \rightarrow 0(\Delta x \rightarrow 0$ 时 $)$
	使得 $|\Delta x|$ 充分小时,
	$$\frac{f\left(x_0+\Delta x\right)-f\left(x_0\right)}{\Delta x}=f^{\prime}\left(x_0\right)+\varphi(\Delta x)$$那么$$\Delta y=f\left(x_0+\Delta x\right)-f\left(x_0\right)=f^{\prime}\left(x_0\right) \Delta x+\varphi(\Delta x) \Delta x=f^{\prime}\left(x_0\right) \Delta x+o(\Delta x)$$
	从而 $y=f(x)$ 在 $x=x_0$ 处可微且 $\mathrm{d}y =f^{\prime}(x) \mathrm{d} x$
	
	充分性：已知 $y=f(x)$ 在 $x=x_0$ 处可微，存在一个仅与 $x_0$ 有关的量 $A(x_0)$使得
	$$\begin{aligned}
		\Delta y=f\left(x_0+\Delta x\right)-f\left(x_0\right)=A\left(x_0\right) & \Delta x+o(\Delta x) \\
		\frac{\Delta y}{\Delta x}=\frac{f\left(x_0+\Delta x\right)-f\left(x_0\right)}{\Delta x} & =A\left(x_0\right)+\frac{o(\Delta x)}{\Delta x} \\
		& \rightarrow A\left(x_0\right)(\Delta x \rightarrow 0 \text { 时 })
	\end{aligned}$$
	$f(x)$ 在 $x=x_0$ 处可导，且 $A\left(x_0\right)=f^{\prime}\left(x_0\right)$ ，故有 $\mathrm{d} y=f^{\prime}(x) \mathrm{d} x$
	
	
\end{proof}
综上 $ f^{\prime}(x)=\dfrac{\mathrm{d} y}{\mathrm{d} x}$

对 $\mathbb{R}^{1}$ 上的函数，可导等价于可微.

对上式的两端,$f^{\prime}(x)$ 是一个整体,$\dfrac{\mathrm{d} y}{\mathrm{d} x}$ 又称微商,它可以做乘除的运算.
\begin{proposition}
	$y=f(x)$ 在 $x=x_0$ 处可导，则 $f(x)$ 在 $x=x_0$ 处连续.
\end{proposition}


\begin{proof}
	设 $f(x)$ 在 $x=x_0$ 处可导，$|\Delta x|$ 充分小时,$\exists \varphi(x) \rightarrow 0(x \rightarrow 0)$,使得$$f\left(x_0+\Delta x\right)-f\left(x_0\right)=f^{\prime}\left(x_0\right) \Delta x+\varphi(\Delta x)$$
	
	$\Delta x \rightarrow 0$ 时，上式右侧$\rightarrow 0$ ,故$$\lim _{\Delta x \rightarrow 0} f\left(x_0+\Delta x\right)=f\left(x_0\right)$$
	
	这就是连续性的定义.
\end{proof}
连续不一定可导，例如 $y=|x|$ 在 $x=0$ 处


